\label{sec:related_work}

\subsection{足式机器人的模型预测控制}

模型预测控制（MPC）已经成为动态足式机器人控制的主流框架之一，因为它能够在统一的滚动优化框架中耦合运动生成、接触调度以及多类约束处理。早期代表性工作表明，即使采用相对简化的预测模型，也能获得很强的动态运动性能，例如 MIT Cheetah 3 上的凸 MPC \cite{dicarlo2018convexmpc}。随后，该方向不断向更丰富的非线性与全身优化建模推进。Grandia 等 \cite{grandia2023perceptive} 证明了非线性 MPC 可以与感知驱动的落脚点规划相结合，实现复杂地形上的动态运动；Meduri 等 \cite{meduri2022biconmp} 提出了能够在线生成周期与非周期动作的全身非线性 MPC 框架。与此同时，多接触规划与全身优化也被进一步拓展到 loco-manipulation 场景 \cite{sleiman2023versatile}，使优化问题的维度和多模态特性进一步增强。

上述工作共同说明，MPC 在足式机器人领域已经具备足够强的表达能力。然而，现有研究更强调地形感知、全身规划和任务泛化，而对于\emph{机构层面的分支一致性}关注相对有限。换言之，已有工作较少讨论这样一个问题：当关节机构存在多种看似相近的解分支时，在线优化器是否能够始终停留在物理上正确的那一支。本文正是在六足机器人背景下研究这一问题，因为更多的腿和更多的接触会显著放大求解器对错误分支局部极小值的敏感性。

\subsection{高维接触优化中的局部极小值陷阱}

足式机器人优化的困难不仅来自状态维数，更来自接触驱动的非凸性和混杂模式切换。Posa 等 \cite{posa2014direct} 指出，接触丰富的轨迹优化本质上是困难的，因为可行运动依赖于不连续的接触事件、互补约束以及模式相关动力学。近期面向敏捷运动的预测控制工作同样承认，在线求解高非线性、多接触优化问题仍然极具挑战 \cite{mastalli2022agile}。在这样的场景下，局部极小值尤其具有欺骗性：两个解在任务空间上可能极为接近，但其对应的关节空间分支并不一定同样可执行、同样物理合理。

这一问题在 warm-start 驱动的 SQP 流程中尤为突出。当局部二次子问题围绕一个已经带偏的初值展开时，若错误分支没有被显式地从可行域中剔除，求解器就可能沿着这一局部吸引盆持续收敛。虽然 runtime guard 可以在后端拦截最终指令，但这只是掩盖了病灶，而没有真正解决它。因此，本文认为，对于多足系统而言，局部极小值问题不应仅被看作优化层面的数值困难，更应被理解为一种与安全性直接相关的建模问题。

\subsection{内点法与硬约束处理}

内点法为约束最优控制提供了另一种思路：它将不等式约束直接纳入原始-对偶求解过程，而不是将其弱化为外部罚函数偏好。经典非线性规划文献已系统证明，带 filter line search 的原始-对偶内点法在大规模约束问题上具有很强的鲁棒性 \cite{wachter2006ipopt}。在嵌入式 MPC 场景中，Frison 和 Diehl \cite{frison2020hpipm} 进一步表明，高性能 IPM 实现可以高效求解结构化最优控制 QP，同时保持良好的约束处理能力。

尽管如此，足式机器人 MPC 中的大量实现仍优先使用 SQP、实时迭代或凸化近似，其原因在于这些方法具有成熟的工程实现与较高的速度优势。本文的结果表明，当控制器必须处理\emph{机构特定的硬不等式约束}时，这种偏好可能变得脆弱。与大多数仅将硬约束用于接触、摩擦锥或落脚点可行性的工作不同，本文引入的 \emph{branch-consistency constraint} 直接对左右 tarsus 的物理解分支进行编码。通过将这一约束与 IPM 后端配对，我们使分支有效性成为求解器内部的可行性条件，而不是 runtime 层的补救规则。

\subsection{本文的定位}

本文与既有足式 MPC 工作的区别体现在“问题诊断”和“修复方式”两个层面。诊断层面，我们揭示了一种隐藏的六足机器人失效模式：优化器可以收敛到任务空间近似合理、但机构分支物理上错误的 tarsus 解。修复层面，我们表明，求解器可见的 \emph{TarsusBranchConstraint} 与 IPM 后端的结合可以在实时控制中消除此类失效。因此，本文的贡献并不仅仅是“又一个非线性 MPC 控制器”，而是一种以可行域塑形为核心的在线优化重构方式，它把物理分支有效性、runtime 安全性和求解器设计统一起来。
